\section{Untying the children introduces a second escape}
\label{sec:untied}

Suppose the two subtrees have independent trainable tables $h_L,h_R$, with the
same grammar, heads, objective and gradient-flow metric as in
Section~\ref{sec:parity-toy}. The easy contrast $U$ is directly available
only on the left. When the left table learns $U$, the right table must
still learn $V$ to fit the interaction. In this model the time spent
waiting for the right table grows logarithmically.

Write $H_L=U,H_R=V$, and for $i=L,R$ put
$a_i=\langle h_i,H_i\rangle$, $R_i=\langle h_i,h_i\rangle$ and
$r_i=a_i/\sqrt{R_i}$. Initialize the tables with $R_L=R_R=1$,
$r_L=r_R=\delta>0$, and $b=w=0$. The nuisance components of the two
tables may have different directions. The hard variance fraction now obeys
\[
 D=r_L^2r_R^2-\big(w\sqrt{R_LR_R}-r_Lr_R\big)^2,
 \qquad \max_w D=r_L^2r_R^2 .
\]
In particular $D\ge p$ implies $r_R\ge\sqrt p$ even if the left table has
already learned its hidden sign perfectly.

\paragraph{Equations.}
By independence, the full population loss is
\[
 2\mathcal L=s^2+1+b^2R_L+w^2R_LR_R-2sba_L-2wa_La_R.
\]
Taking its derivatives gives
\begin{align*}
 \dot b&=sa_L-bR_L,& \dot w&=a_La_R-wR_LR_R,\\
 \dot h_L&=(sb+wa_R)H_L-(b^2+w^2R_R)h_L,&
 \dot h_R&=wa_L H_R-w^2R_Lh_R .
\end{align*}
The inner products of the last two equations with $2h_L,2h_R$ give
$\dot R_L=2b\dot b+2w\dot w$ and $\dot R_R=2w\dot w$.
Projecting each table velocity perpendicular to the current table gives
the alignment equations. Together these yield
\begin{align}
 R_L&=1+b^2+w^2, & R_R&=1+w^2,\nonumber\\
 \dot r_L&=(1-r_L^2)(sb+wa_R)/\sqrt{R_L},&
 \dot r_R&=(1-r_R^2)w a_L/\sqrt{R_R}.
 \label{eq:untied-flow}
\end{align}

\begin{theorem}
\label{thm:untied}
Fix $q_L,q_R,p\in(0,1)$. Let $T_L(q_L),T_R(q_R),T_p(p)$ denote the
first times that the left alignment reaches $q_L$, the right alignment
reaches $q_R$, and the trained root explains $p$ of the interaction
variance, respectively. All three thresholds are attained for sufficiently
small $\delta>0$. If $s=1$, then, with constants in the remainders
depending on the fixed thresholds,
\begin{align*}
 T_L(q_L)&=\lambda_1^{-1}\log(1/\delta)+O(1),\\
 T_R(q_R),\ T_p(p)&=(\lambda_1^{-1}+\lambda_2^{-1})\log(1/\delta)+O(1),
\end{align*}
where $\lambda_1=(\sqrt5-1)/2$,
$\phi=(1+\sqrt5)/2$, and
$\lambda_2=(\sqrt{\phi^2+4\phi}-\phi)/2$.
If $s=0$, each of the three times satisfies
\[
 T=\frac{1}{2\delta^2}+3\log(1/\delta)+O(1).
\]
\end{theorem}

\begin{proof}
The readout equations have equilibrium targets $sr_L/\sqrt{R_L}$
and $r_Lr_R/\sqrt{R_LR_R}$. Both lie in $[0,1]$ on the positive unit
cube. The alignment equations preserve the same cube, with both alignments
nondecreasing, so $1\le R_L\le3$ and $1\le R_R\le2$. The solution exists
globally. Integrating the damped readout equations from zero and using
monotonicity yields the sharper bounds $b\le r_L$ and $w\le r_Lr_R$.

\emph{The left table escapes first when $s=1$.}
Let $t_1$ be the first time $r_L=\varepsilon$, for a small fixed
$\varepsilon<q_L$, and put $x=r_L+\lambda_1 b$. Before $t_1$,
$R_L=1+O(r_L^2)$, $R_R=1+O(r_L^2)$ and $r_L\le x\le(1+\lambda_1)r_L$.
The $b$ equation gives $\dot b=r_L-b+O(r_L^3)$, and the positive
interaction term in $\dot r_L$ gives $\dot r_L\ge b-O(r_L^3)$.
Using $\lambda_1^2+\lambda_1=1$, we obtain
$\dot x\ge\lambda_1x-O(x^3)\ge\lambda_1x/2$. This also proves
that $t_1$ is finite. Hence $\int_0^{t_1}r_L^2\,dt$ is bounded
independently of $\delta$, by integrating $x^2$ against
$dt\le2\,dx/(\lambda_1x)$.
The second alignment equation and $w\le r_Lr_R$ now give
\[
 0\le\frac{d}{dt}\log r_R\le C r_L^2 .
\]
Thus $\delta\le r_R\le C\delta$ until $t_1$. It follows that
$wa_R=O(r_L\delta^2)=O(r_L^3)$ on this interval, so the lower estimate
for $\dot x$ sharpens to $\dot x=\lambda_1 x+O(x^3)$.
Integration from $x(0)=\delta$ gives
$t_1=\lambda_1^{-1}\log(1/\delta)+O(1)$.

\emph{The right table escapes with a different exponent.}
Measure time $t$ from $t_1$. We start with $r_R$ between $\delta$ and
$C\delta$, and $0\le w\le r_R$. Since $r_L\ge\varepsilon$, the target
of the $b$ equation stays above $\varepsilon/\sqrt3$. Comparison shows
that $b$ is bounded below by a fixed positive constant after one time
unit. Consequently $1-r_L\le Ce^{-\kappa t}$ for some $\kappa>0$, by
the first alignment equation. In particular the left alignment passes
any fixed $q_L<1$ within bounded time from $t_1$.

With $r_L=1,w=0$, the easy readout would obey
$\dot b=F(b)=\sqrt{1+b^2}-b(1+b^2)$.
Its unique positive fixed point is $b_*=1/\sqrt\phi$, at which
$R_L=\phi$ and $a_L=\sqrt\phi$. Its derivative
$F'(b)=b/\sqrt{1+b^2}-1-3b^2$ is bounded strictly below zero on $[0,1]$.
In the actual system the perturbation to $F(b)$ is
$O(1-r_L+w^2)=O(e^{-\kappa t}+r_R^2)$.
Scalar variation of constants, using monotonicity of $r_R$, therefore gives
\[
 |b-b_*|+|a_L-\sqrt\phi|+|R_L-\phi|
 \le C(e^{-\kappa' t}+r_R^2)
\]
for some $\kappa'>0$. A decrease in the exponent absorbs the convolution
of two decaying exponentials.

While $r_R$ stays small, set $z=(r_R,w)^\top$. Since $a_R=r_R\sqrt{R_R}$
and $R_R=1+O(r_R^2)$, the exact equations take the form
\[
 \dot z=
 \begin{pmatrix}0&\sqrt\phi\\ \sqrt\phi&-\phi\end{pmatrix}z
 +O\big((e^{-\kappa't}+r_R^2)\|z\|\big).
\]
The positive eigenvalue of this matrix is $\lambda_2$ in the theorem.
For its unstable coordinate $y=r_R+(\lambda_2/\sqrt\phi)w$,
positivity and $w\le r_R$ give $r_R\le y\le C r_R$ and
\[
 \frac{d}{dt}\log y=\lambda_2+O(e^{-\kappa't}+r_R^2).
\]
Choose a small fixed $\eta<\min(q_R,\sqrt p)$.
The right alignment remains of order $\delta$ during any fixed initial
interval after $t_1$, by $d\log r_R/dt\le C r_L^2\le C$. After a fixed
transient, the last display gives $\dot y\ge\lambda_2 y/2$ until
$r_R=\eta$. Its hitting time is finite, and $\int r_R^2\,dt$ over
this interval is uniformly bounded using the same change of variable
as for $x$. Thus the integral of the logarithmic-rate error is bounded.
Since $y$ started at order $\delta$, the additional time to reach
$r_R=\eta$ is $\lambda_2^{-1}\log(1/\delta)+O(1)$.

\emph{Fitting the root after both escapes.}
From $r_L\ge\varepsilon$, $r_R\ge\eta$, the target
$r_Lr_R/\sqrt{R_LR_R}$ of the damped $w$ equation is bounded away from
zero, so $w$ becomes bounded below in bounded time. The two positive
alignment equations then give $r_L,r_R\to1$. At any accumulation point,
the squared-loss gradient vanishes by boundedness and dissipation, as in
Theorem~\ref{prop:toy-gap}. Thus $w\sqrt{R_LR_R}=1$ there and $D\to1$.
The compactness and continuous-dependence argument in that theorem gives
uniformly bounded additional times to each fixed $q_R,p<1$. The choice
$\eta<\sqrt p$ and the frozen-child bound exclude an earlier hard
crossing. This proves the two logarithmic estimates.

\emph{Hard contrast alone.}
For $s=0$, $b=0$ and identical initial alignments and norms give
$r_L=r_R=r$ for all time by uniqueness. Set $u=r^2$, $R=1+w^2$ and
$k=w/u$. Equation~\eqref{eq:untied-flow} and the root equation give
\[
 \dot u=2wu(1-u),\quad
 \dot k=R-\big[R^2+2w(1-u)\big]k,\quad
 \frac{d}{dt}u^{-1}=-2k(1-u).
\]
Choose $u_1>0$ small and below $q_L^2,q_R^2,\sqrt p$. It is reached
for every sufficiently small $\delta>0$, since $r\ge\delta$ and the
root equation eventually bounds $w$ below by a positive constant.
On the interval up to this hitting time $t_1$, we have $0\le w\le u$.
Use $g(u)=u-2u^2$, put $e=w-g(u)$, and write its exact equation as
\[
 \dot e=B(w,u),\qquad B(w,u)=R(u-Rw)-2g'(u)wu(1-u).
\]
Here $B(g(u),u)=O(u^3)$ and $\partial_wB=-1+O(u)$ uniformly for
$0\le w\le u\le u_1$. Using $e(0)=-g(\delta^2)$ and monotonicity of $u$
in variation of constants gives
$|e(t)|\le C\delta^2 e^{-t/2}+Cu(t)^3$.
With $F(u)=-1/(2u)+(3/2)\log u$ we therefore have
\[
 \frac{d}{dt}F(u(t))=(1+3u)(1-u)\frac{w}{u}
 =1+O(u^2)+O(e^{-t/2}).
\]
The integral of $u^2$ to $t_1$ is bounded uniformly in $\delta$.
Indeed $\dot k\ge1-C_1k$ bounds $k$ below by a positive constant
after one time unit, and $\dot u=2ku^2(1-u)$ bounds the remaining integral.
It follows that $t_1=F(u_1)-F(\delta^2)+O(1)$. The uniformly bounded
remaining time to any of the three thresholds gives the stated expansion.
\end{proof}

Thus releasing the tying constraint preserves the speedup from adding
the easy contrast, but changes which interval is long. In the tied model
the right branch receives the left branch's feature as soon as the shared
table rotates. In the untied model the right alignment is still of order
$\delta$ when the left alignment first reaches a small fixed level.
Consequently the gap from a fixed left threshold to a fixed hard fit
itself grows as $\lambda_2^{-1}\log(1/\delta)+O(1)$.
